\subsection{Graph of a Function} 
We can recognize the graph of a function.
\begin{itemize}
\item To be a function, we must have each \makeblanksmall-coordinate occur only once.
\item On a graph, the points with the same \makeblanksmall-coordinate form a \makeblank[1in] line.
\item So in the graph of a function, each \makeblank crosses the graph \makeblank.
\item This is called the \makeblank.
\end{itemize}
\begin{example}
Which of the relations graphed below is a function?
\begin{lpic}{}{12}
\setgraphsquare{5}
\begin{scope}[xshift=3\linespace,yshift=-3\linespace,scale=0.5]
\AxesGen{\small}{\tiny}
\draw[graphend] (1,1) circle ({sqrt(10)});
\draw[graphend] (5,5) node[anchor=north east]{$R$};
\end{scope}
\begin{scope}[xshift=3\linespace,yshift=-9\linespace,scale=0.5]
\AxesGen{\small}{\tiny}
\draw[graph,domain=-3:3] plot (\x,\x*\x*\x/3-4*\x/3);
\draw[graphend] (5,5) node[anchor=north east]{$S$};
\end{scope}
%\begin{scope}[xshift=3\linespace,yshift=-15\linespace,scale=0.5]
%\AxesGen{\small}{\tiny}
%\end{scope}
\end{lpic}
\end{example}
%We have two ways to describe the graph of a function:
%\begin{itemize}
%\item The graph of the relation $f$
%\item The graph of the equation $y=f(x)$. This means that, for each point $(x,y)$ on the graph, \makeblank.
%\end{itemize}
